\scanpage{14}
donc
\[
\begin{aligned}
\varepsilon_0^2&=l_0\\
\varepsilon_1^2&=l_1\\
\varepsilon_\infty^2&=l_\infty
\end{aligned}
\]
Ce qui redonne les $l_i$, et aussi les $\sigma_i$, \deleted{par} en termes des $\varepsilon_i$
\[
\begin{aligned}
\sigma_0&=\varepsilon_0\varepsilon_1^2\\
\sigma_1&=\varepsilon_1\varepsilon_\infty^2\\
\sigma_\infty&=\varepsilon_\infty\varepsilon_0^2.
\end{aligned}
\]
D'où les relations
\[
\left\{\begin{aligned}
\varepsilon_0^2\varepsilon_1^2\varepsilon_\infty^2&=1\\
(\varepsilon_0\varepsilon_1^2)^2&=1\\
(\varepsilon_1\varepsilon_\infty^2)^2&=1\\
(\varepsilon_\infty\varepsilon_0^2)^2&=1
\end{aligned}\right.
\]
On retrouve $\rho$ \deleted{comme} par
\deleted{\[
\rho=\sigma_0\sigma_1=\sigma_1\sigma_\infty=\sigma_\infty\sigma_0
\]}
\deleted{\[
\begin{gathered}
\Vert\\
\varepsilon_0\varepsilon_1^3\unsure{\varepsilon_\infty|\varepsilon_1}
\end{gathered}
\]}
\inserted{\[
\sigma_0\sigma_1=l_1^{-1}\rho
\]}
\margin{donc
\[
\begin{aligned}
\rho&=l_1\sigma_0\sigma_1\\
&=\varepsilon_1^2\varepsilon_0\varepsilon_1^2\varepsilon_1\varepsilon_\infty^2\\
&=\varepsilon_1^2\varepsilon_0\varepsilon_1^{\deleted{\unsure{2}}}\underbrace{\varepsilon_1^2\varepsilon_\infty^2}_{\varepsilon_0^{-2}}
\end{aligned}
\]}
\[
\begin{aligned}
\rho&=\varepsilon_1^2\varepsilon_0\varepsilon_1\varepsilon_0^{-2}\\
&=\varepsilon_\infty^2\varepsilon_1\varepsilon_\infty\varepsilon_1^{-2}\\
&=\varepsilon_0^2\varepsilon_\infty\varepsilon_0\varepsilon_\infty^{-2}.
\end{aligned}
\]
\margin{ces relations ne semblent pas \unsure{conséquences} des autres}
Multipliant \deleted{dans} ces relations dans l'ordre inverse et écrivant $\rho^3=1$, on trouve
\[
(\varepsilon_0\varepsilon_1\varepsilon_\infty)^2=1
\]
ce qui est inquiétant, car il n'y a pas d'\unsure{élém\textsuperscript{ts}}
\mspage{13}
